%% -*- coding: utf-8; -*-
%% Dissertação de Mestrado — PUC-Rio (classe oficial thesispuc)
%% Compilar: pdflatex dissertacao && bibtex dissertacao && pdflatex dissertacao && pdflatex dissertacao
%% Corpo em corpo.tex; bibliografia em referencias.bib.
%% Campos marcados com % TODO precisam de confirmação antes da versão final.
\documentclass[msc,brazilian,digital]{thesispuc}

%%% Pacotes adicionais
\usepackage{tabularx}
\usepackage{multirow}
\usepackage{multicol}
\usepackage{colortbl}
\usepackage[dvipsnames,svgnames,x11names,fixpdftex,table]{xcolor}
\usepackage{numprint}
\usepackage{textcomp}
\usepackage{booktabs}
\usepackage{amsmath}
\usepackage{amssymb}
\usepackage{enumitem}
\usepackage[alf,bibjustif,abnt-emphasize=bf]{abntex2cite}

%% --- Legendas de tabelas e figuras -----------------------------------
%% Pedido recorrente da revisao: a legenda vinha no mesmo corpo e na mesma
%% fonte do texto e se confundia com ele. Aqui ela fica menor, com o rotulo
%% em negrito, travessao no lugar dos dois-pontos e recuo das margens do
%% bloco de texto.
\captionsetup{font=footnotesize,labelfont=bf,labelsep=endash,
              justification=raggedright,singlelinecheck=false,
              margin=14pt,skip=6pt,belowskip=2pt}

%% --- Colocacao de floats ---------------------------------------------
%% Sem afrouxar estes parametros o LaTeX empurra tabela grande para pagina
%% propria, longe do trecho que a discute.
\renewcommand{\topfraction}{0.85}
\renewcommand{\bottomfraction}{0.70}
\renewcommand{\textfraction}{0.10}
\renewcommand{\floatpagefraction}{0.75}
\setcounter{topnumber}{3}
\setcounter{totalnumber}{4}

\npthousandsep{.}
\npdecimalsign{,}

\graphicspath{{figuras/}{./}}

\setcounter{tocdepth}{1}
\usecolour{true}

%% Seções opcionais da classe desativadas (esta dissertação não as usa)
\abreviationsmode{none}
\codesmode{none}
\algorithmsmode{none}

%%% Autoria
\author{Mariana Barreto}
\authorR{Barreto, Mariana}

\advisor{Sérgio Lifschitz}{Prof.}
\advisorR{Lifschitz, Sérgio}

\title{O Efeito do Volume de Coleta sobre os Resultados de Análises: Redes Sociais Digitais como Caso de Estudo}
\titleuk{The Effect of Collection Volume on the Results of Analyses: Online Social Networks as a Case Study}

\day{29}          % data da defesa (Formulário de Marcação de Defesa)
\month{Setembro}
\year{2026}

\city{Rio de Janeiro}
\CDD{004}
\department{Informática}
\program{Informática}
\school{Pós-Graduação em Informática}
\university{Pontifícia Universidade Católica do Rio de Janeiro}
\uni{PUC-Rio }

%%% Banca (a classe acrescenta o orientador; aqui vão os demais membros)
\jury{%
  \jurymember{Edward Hermann Haeusler}{Prof.}
    {Departamento de Informática}{PUC-Rio}
  \jurymember{Ana Carolina Brito de Almeida}{Profa.}
    {Instituto de Matemática e Estatística}{UERJ}
}

%%% Resumo pessoal (bio)
\resume{%
Bacharel em Engenharia da Computação pela Pontifícia Universidade Católica do
Rio de Janeiro (PUC-Rio) em 2023. Ingressou no mestrado em Informática da mesma
instituição em 2024, com pesquisa em bancos de dados e em análise de dados de
redes sociais digitais.%
}

\acknowledgment{%
\noindent Ao meu orientador, Professor Sérgio Lifschitz, pelo estímulo e
parceria ao longo deste trabalho.
\bigskip

\noindent Aos professores Edward Hermann Haeusler e Ana Carolina Brito de
Almeida, por aceitarem compor a banca examinadora e pela leitura atenta deste
trabalho.
\bigskip

\noindent À minha família, pelo apoio e pela paciência ao longo de todo o
mestrado.
\bigskip

\noindent O presente trabalho foi realizado com apoio da Coordenação de
Aperfeiçoamento de Pessoal de Nível Superior --- Brasil (CAPES) --- Código de
Financiamento 001.
}

\catalogprekeywords{%
  \catalogprekey{Informática}%
}

\keywords{%
  \key{Volume de Coleta de Dados}
  \key{Redes Sociais Digitais}
  \key{Amostragem}
  \key{Análise de Redes Sociais (SNA)}
  \key{Validade de Conclusões}
  \key{Reprodutibilidade}
}

\keywordsuk{%
  \key{Data Collection Volume}
  \key{Online Social Networks}
  \key{Sampling}
  \key{Social Network Analysis (SNA)}
  \key{Validity of Conclusions}
  \key{Reproducibility}
}

\abstract{%
Para as análises mais realizadas sobre Redes Sociais Digitais (RSD), como detecção
de comunidades, análise de conteúdo, sentimento e classificação, coletar pouco ou
muito muda a conclusão? A pergunta é uma instância de um problema metodológico
mais amplo, o do efeito do volume de coleta sobre os resultados de uma análise.
Esta dissertação toma as RSD como caso de estudo, por serem um domínio em que esse
efeito é mensurável. Um levantamento sistemático feito como parte do
projeto~\cite{barreto2026survey} mostra que o volume coletado varia enormemente,
mas que menos da metade dos artigos passa de 1 milhão de publicações coletadas,
e que quase nunca se avalia se esse volume altera a conclusão. Para medir o efeito, propõe-se um protocolo que trata cada
artigo como uma amostra sub-coletada do seu tema, coleta um superconjunto temático
e reproduz a análise sobre a coleta original e a expandida, testando se a conclusão
muda e se já havia convergido (caso em que coletar menos bastaria). Um conjunto de
casos replicáveis, distribuídos por diferentes redes sociais e tipos de análise,
instancia o protocolo. A entrega não é um volume ideal único, mas a identificação
de para quais tipos de análise o volume de coleta é decisivo e para quais é
indiferente, de modo que declarar o volume coletado, e o quanto a conclusão depende
dele, passe a integrar o relato de qualquer trabalho que colete dados de RSD.%
}

\abstractuk{%
For the analyses most frequently performed over Online Social Networks (OSN), such as
community detection, content analysis, sentiment and classification, does collecting
little or a lot change the conclusion? The question is an instance of a broader
methodological problem, that of the effect of collection volume on the results of
an analysis. This dissertation takes OSN as a case study, as a domain where the
effect is measurable. A systematic survey conducted as part of the
project~\cite{barreto2026survey} shows that the collected volume varies enormously,
yet that fewer than half of the papers exceed one million collected posts, and
that whether that volume changes the conclusion is almost never assessed. To measure the effect, we propose a
protocol that treats each paper as an under-collected sample of its theme, collects
an on-theme superset and reproduces the analysis over the original and the expanded
collection, testing whether the conclusion changes and whether it had already
converged (in which case less would suffice). A lineup of replicable cases, spread
across different social networks and analysis types, instantiates the protocol. The
deliverable is not a single ideal volume but an account of the analysis types for
which collection volume is decisive and those for which it is immaterial, so that
stating the collected volume, and how much the conclusion depends on it, becomes
part of the report of any work that collects OSN data.%
}

\dedication{%
  Aos meus pais, Mauro e Regina, e ao meu irmão Marcelo, que estiveram do meu
  lado em todas as minhas jornadas.
}

%%% Epígrafe: opcional, ainda não escolhida. As três linhas ficam comentadas
%%% de propósito. A classe liga a página de epígrafe ao ver a chamada de
%%% \epigraph, e não pelo conteúdo dela, de modo que chamá-la vazia
%%% imprime uma página em branco com a vírgula e o ponto do crédito.
% \epigraph{%
%   TODO: epígrafe (opcional)
% }
% \epigraphauthor{}
% \epigraphbook{}

\begin{document}
  % ============================================================
% A Introdução (Cap. 1) será escrita ao final e está em
% introducao.tex, NÃO incluída no PDF. O capítulo abaixo fica FORA da
% introdução e é compilado.
% ============================================================

% =====================================================================
\chapter{Contexto, Objetivos e Contribuições}
\label{cap:objetivos}
% =====================================================================

\section{Contexto: o levantamento realizado}
\label{sec:contexto-survey}

Para fundamentar empiricamente a motivação, conduziu-se um levantamento
sistemático da literatura que emprega coleta automatizada em RSD, reportado em
artigo autônomo~\cite{barreto2026survey}. Em resumo, um fluxo reprodutível partiu
do índice DBLP, deduplicou e enriqueceu os metadados, baixou os PDFs em acesso
aberto e extraiu informações estruturadas sobre a coleta por meio de três modelos
de linguagem independentes, com validação contra gabarito manual. Foram
catalogados 13.395 artigos; 2.139 tiveram a coleta caracterizada por extração
automática; e, após excluir 421 que, apesar de extraídos, não coletam dados de RSD
(entrevistas e estudos qualitativos), restaram 1.718 artigos como base de análise
da coleta.

Dois achados desse levantamento motivam diretamente esta dissertação. Primeiro, a
distribuição dos tipos de análise entre os 2.139 artigos extraídos
(Tabela~\ref{tab:analises-survey}), que orienta a escolha dos métodos a investigar.
Segundo, e central para este trabalho, entre os 1.718 que efetivamente coletam RSD
o volume coletado varia por mais de sete ordens de grandeza (mediana de
$\approx$273 mil itens, mas 40\% dos artigos passam de 1 milhão e 22\% coletam ao
menos 10 milhões), e coletar muito não implica amostrar. Entre os maiores
coletores, quase nenhum amostra estatisticamente, e a fração de fato analisada
pode ser ínfima; um estudo coletou $\approx$2,8 bilhões de tweets e analisou
apenas $\approx$1\%. A literatura dimensiona o volume por conveniência e
tipicamente não verifica se esse recorte preserva as conclusões, e é esse vão que
o protocolo desta dissertação mede.

\begin{table}[t]
\centering
\caption{Tipos de análise no corpus do levantamento ($N=2.139$
artigos)~\cite{barreto2026survey}. ``Análise de conteúdo'' (1.396 menções) é
omitida por ser um guarda-chuva, não um método.}
\label{tab:analises-survey}
\begin{tabular}{lrr@{\hspace{2em}}lrr}
\toprule
Tipo & Art. & \% & Tipo & Art. & \% \\ \midrule
Estatística descritiva & 935 & 43,7 & Modelagem de tópico   & 408 & 19,1 \\
Sentimento             & 865 & 40,4 & Detecção de comunidades & 334 & 15,6 \\
Classificação          & 756 & 35,3 & Centralidade          & 319 & 14,9 \\
Temporal               & 665 & 31,1 & Influência            & 237 & 11,1 \\
Análise de rede        & 527 & 24,6 & Clusterização        & 157 & 7,3 \\
\bottomrule
\end{tabular}
\end{table}

\section{Objetivos}

O objetivo central é propor e instrumentar um protocolo que meça, para os tipos de
análise mais recorrentes em RSD, o efeito do volume de coleta sobre a conclusão.
Cada artigo é tratado como uma amostra sub-coletada do seu tema; coleta-se um
superconjunto temático e verifica-se se coletar mais muda a conclusão e se ela já
havia convergido no volume original, caso em que coletar menos teria bastado. Mais
do que caracterizar as RSD, busca-se produzir, no domínio em que o efeito é
mensurável, evidência transferível sobre para quais tipos de análise o volume de
coleta é decisivo, e não uma prescrição numérica única de quanto coletar. O
levantamento (\S\ref{sec:contexto-survey}) é o objetivo instrumental que fundamenta
este.

\section{Perguntas de pesquisa}

\begin{description}[leftmargin=2.2em,style=nextline]
  \item[PP1] Para cada tipo de análise, coletar mais do que o artigo coletou (um
  superconjunto temático) muda a conclusão? Isto é, coletar mais era necessário?
  \item[PP2] A conclusão já havia convergido no volume originalmente coletado, caso
  em que esse volume bastava, ou até excedia o necessário, ou ainda se movia,
  indicando coleta prematura?
  \item[PP3] Para quais tipos de análise o volume de coleta é uma variável crítica,
  com a conclusão dependendo de se coletar mais ou menos, e para quais é
  praticamente inócuo?
\end{description}

\section{Contribuições}

\begin{itemize}
  \item um protocolo reprodutível para medir o efeito do volume de coleta, por tipo
  de análise, com testes de mudança e de convergência da conclusão, isto é, se
  coletar mais era necessário e se coletar menos teria bastado, acrescido de um
  teste que separa falha de volume de falha de esquema de amostragem
  (Seção~\ref{sec:curva-midia});
  \item um conjunto fixado de seis casos replicáveis
  (Capítulo~\ref{cap:experimento}), um por rede social, com a expansão de coleta
  especificada e ainda acessível hoje, dois deles já executados
  (Capítulos~\ref{cap:caso-twitter} e~\ref{cap:caso-vacinas});
  \item como contribuição instrumental, o levantamento sistemático e seu método de
  extração multi-modelo, reportados em artigo autônomo~\cite{barreto2026survey};
  \item enquanto caso de estudo, evidência concreta, num domínio em que o efeito é
  mensurável, de para quais análises o volume de coleta muda a conclusão,
  informando um problema metodológico que extrapola as RSD.
\end{itemize}

% =====================================================================
\chapter{Metodologia}
\label{cap:experimento}
% =====================================================================

O levantamento estabelece o quê. A literatura quase sempre dimensiona a coleta
por conveniência e raramente justifica que o volume coletado preserva as
conclusões. Este capítulo trata de como medir o efeito do volume (se coletar mais
era necessário e se coletar menos bastaria), por tipo de análise. Tanto encontrar
que a conclusão muda quanto encontrar que ela é robusta ao volume são resultados
úteis.

\section{Lógica do protocolo}

Para cada artigo-alvo, reproduz-se o fluxo de análise
tal-como-publicado sobre a coleta original do artigo ($A_{\text{paper}}$);
coleta-se um superconjunto temático, com mais dados sobre o mesmo tema, e roda-se
a mesma análise sobre ele ($A_{\text{mais}}$); constrói-se a curva
$A(\text{volume})$ entre os dois pontos. A decisão usa dois testes. No primeiro,
verifica-se se $A_{\text{paper}}$ difere de $A_{\text{mais}}$, com divergência
acima da banda bootstrap e inversão da afirmação do artigo. No segundo, se a curva
ainda se movia em $N_{\text{paper}}$, a coleta do artigo era prematura. O critério
de diferença material é fixado antes da execução, em duas camadas: a estatística,
em que o resultado sai do intervalo de confiança bootstrap; e a substantiva, em
que a afirmação do artigo se inverte (por exemplo, ``o usuário $X$ é o mais
central''). A camada substantiva é a decisiva. A combinação dos dois testes evita
declarar erro apenas por haver diferença: como a expansão é temática, ela aproxima
a população do tema, e a não-convergência mostra que a conclusão não estava
estabilizada.

\section{Eixos de expansão (mantendo o tema)}

Cada artigo ocupa um ponto restrito em um ou mais eixos de coleta; a
expansão percorre esses eixos sem sair do tema. O tema é definido pelo próprio
escopo do artigo (subreddit, hashtags, termos eleitorais), e é esse filtro
temático que se mantém ao coletar mais (Tabela~\ref{tab:eixos}). Constrói-se a curva $A(\text{volume})$ por frações
crescentes, de $N_{\text{paper}}$ a $N_{\text{máx}}$, com ao menos 30 réplicas
bootstrap por ponto quando a coleta permitir subconjuntos.

\begin{table}[t]
\centering
\caption{Eixos de expansão da coleta, mantendo o tema do artigo.}
\label{tab:eixos}
\begin{tabular}{lp{4.0cm}p{4.0cm}}
\toprule
Eixo & Restrição típica do artigo & Expansão temática \\
\midrule
Volume / top-$k$    & ``1.506 publicações do topo'', aba de mais votados & todos os itens do escopo \\
Janela temporal     & ``2 semanas'', ``$\pm$6 meses''      & estender o período \\
Cobertura de fontes & ``4 subreddits'', 100 de 1.973 sementes & todas as fontes do tema \\
Largura do filtro   & poucas hashtags-semente      & ampliar hashtags/termos \\
Granularidade       & só publicações/vídeos                       & incluir comentários/respostas \\
\bottomrule
\end{tabular}
\end{table}

\section{Métricas de divergência por tipo de análise}

A escolha dos tipos a cobrir segue o levantamento (Tabela~\ref{tab:analises-survey}), decompondo ``análise de conteúdo'' em
métodos concretos e privilegiando os de maior sensibilidade à coleta. Cobrem-se
quatro métodos núcleo (com classificação como opcional):

\begin{description}[leftmargin=1.5em,style=nextline]
  \item[Detecção de comunidades] Normalized Mutual Information
  (NMI)~\cite{danon2005nmi} e Adjusted Rand Index (ARI) entre a
  partição da coleta original e a da expandida; variação da modularidade e do
  número de comunidades; sobrevivência das $k$ maiores. Algoritmos Louvain/Leiden
  com a mesma resolução~\cite{blondel2008louvain,traag2019leiden}.
  \item[Centralidade] correlação de ordenação (Spearman, Kendall-$\tau$)
  entre os vetores de centralidade; sobreposição dos top-$k$ (Jaccard,
  rank-biased overlap)~\cite{webber2010rbo}; permanência dos atores
  que o artigo aponta como centrais.
  \item[Modelagem de tópico] alinhamento de tópicos original$\leftrightarrow$
  expandida por similaridade das distribuições palavra-tópico (atribuição
  húngara); número de tópicos que sobrevivem; coerência. Parâmetros e
  vocabulário fixados ao do artigo (LDA/BERTopic~\cite{grootendorst2022bertopic}).
  \item[Análise de sentimento] divergência de Jensen--Shannon entre as
  distribuições de polaridade; deslocamento da classe majoritária; correlação
  das séries temporais de sentimento. Modelo idêntico ao do artigo
  (p.ex.\ VADER~\cite{hutto2014vader}), para não confundir efeito de coleta com
  efeito de modelo.
\end{description}

\section{Coleta expandida por plataforma}

Coletar mais do que o artigo depende da plataforma e dos acessos disponíveis. O
Reddit é o caso ideal. Os arquivos históricos do
Pushshift~\cite{baumgartner2020pushshift} (e sucessores), distribuídos via
Academic Torrents, fornecem a população completa de um
subreddit/janela, enquanto a API oficial (PRAW) entrega apenas
$\sim$1000 itens por listagem; a própria API é, portanto, um esquema de
amostragem a superar. O YouTube (Data API v3) tem amostragem forçada
pelo teto de $\sim$500 resultados por consulta. A Meta Content Library
(Instagram/Facebook/Threads) e a TikTok Research API, acessíveis
institucionalmente, têm cotas (p.ex.\ 500k publicações por semana) que delimitam a
coleta expandida.

\section{Conjunto de casos}
\label{sec:casos}

Os artigos foram selecionados do corpus do levantamento para cobrir os tipos de
análise mais recorrentes numa diversidade de redes sociais, priorizando casos
que sub-coletaram de forma expansível e cuja coleta segue acessível hoje. A
Tabela~\ref{tab:casos} resume o conjunto de casos; em todos, a expansão é temática
(mantém o assunto do artigo). O caso de centralidade (Buntain e Golbeck, 2014) é
exemplar: por limite da API, o artigo coletou apenas os cem posts mais votados do
mês e cerca de duzentos comentários de cada, e ainda descartou usuários com menos
de vinte conexões; como o tema são os treze subreddits analisados, coletar todas
as publicações e comentários é expandir sem sair do tema.

Um sexto caso cobre o Twitter/X, ausente dos demais por já não permitir coleta
nova, mas responsável por 54,5\% do levantamento DBLP e, portanto, uma lacuna
importante. Ele entra por exceção: o laboratório eTC preserva uma
coleta histórica das eleições de 2014 (10,7 milhões de tweets),
o que dispensa coleta nova. É também o primeiro caso levado à execução
(Capítulo~\ref{cap:caso-twitter}). Pela mesma exceção entrou, e também já foi
executado, um segundo caso em Twitter/X, o do debate vacinal de 2021--22
(Capítulo~\ref{cap:caso-vacinas}), igualmente preservado pelo eTC. Ele não
acrescenta uma rede ao conjunto, e sim dois tipos de análise ainda descobertos,
detecção de comunidades e análise de conteúdo, sobre uma sub-coleta de natureza
distinta: não um volume pequeno, mas um estrato recortado por limiar de
engajamento dentro de uma coleta larga. O CrowdTangle, sem coleta nem acervo
equivalente, fica de fora. O caso de Instagram/Facebook tem alvo definido
(Capozzi et al., 2021), e nele a sub-coleta é declarada pelos próprios autores,
que buscam anúncios exibidos apenas no Facebook, não no Instagram.

\begin{table}[t]
\centering\small
\caption{Os seis casos do estudo, cobrindo os tipos de análise mais recorrentes.}
\label{tab:casos}
\begin{tabular}{llp{3.2cm}p{3.9cm}}
\toprule
Análise & Rede (\#) & O que coletou & Coletar mais (no tema) \\
\midrule
Centralidade   & Reddit (Buntain 2014)  & top-100 posts + $\sim$200 coment.\ (jul/2013); corte $\geq$20 arestas ($n{=}279$) & todas as submissões e comentários dos 13 subreddits, sem o corte \\
Mod.\ de tópico & YouTube (AGECovP 2025) & 85 buscas; filtro exige a palavra no título ou na descrição (descarta 52\% da busca e 99\% dos sugeridos) & as mesmas buscas sem o filtro de palavra-chave \\
Análise de conteúdo & TikTok (PoliTok-DE 2025) & 939 mil posts de duas eleições alemãs; disponibilidade checada em datas fixas & eixo temporal: rechecagem em datas sucessivas \\
Classificação & Instagram/Facebook (Capozzi 2021) & anúncios exibidos só no Facebook, \emph{não} no Instagram (declarado pelos autores); 2.312 anúncios & incluir o Instagram; braço Brasil \\
Polarização    & Reddit (Massachs 2020) & focus group de 44.924 usuários ($\geq$10 coment.\ em 2012 e 2016) & relaxar o limiar de atividade (usuários menos ativos) \\
Mídia + stance & Twitter (\#Eleições2014) & 700 tweets (100/dia; 1 hashtag) & universo do acervo eTC (10,7M; + hashtags eleitorais) \\
\bottomrule
\end{tabular}
\end{table}

\section{Protocolo operacional}

O protocolo procede em cinco fases: (0) reprodução do artigo
tal-como-publicado sobre a coleta original, com $A_{\text{paper}}$; (1) coleta de
um superconjunto temático e seu congelamento em instantâneo
versionado, com $A_{\text{mais}}$; (2) construção da curva $A(\text{volume})$ em
frações crescentes de $N_{\text{paper}}$ a $N_{\text{máx}}$; (3) cálculo da
divergência com banda bootstrap e os dois testes (a conclusão
mudou? e já havia convergido em $N_{\text{paper}}$?); (4) síntese por
caso. O entregável central é uma tabela mestre (por análise e rede: a coleta do artigo
bastava?, a conclusão mudou?, fração mínima do tema que a estabiliza), acompanhada
das curvas $A(\text{volume})$, e
uma leitura, por tipo de análise, de quando o volume de coleta é decisivo
para a conclusão.

% =====================================================================
\chapter{Primeiro Caso Executado: \#Eleições2014 no Twitter}
\label{cap:caso-twitter}
% =====================================================================

Este capítulo relata o primeiro caso do conjunto de casos levado à execução: a
replicação por expansão de \citeonline{ituassu2015temas}, que analisou o público
da hashtag \#Eleições2014 no segundo turno das eleições presidenciais
brasileiras de 2014. É o sexto caso (\S\ref{sec:casos}) e o único em Twitter/X. Embora a coleta nova nessa plataforma seja hoje inviável, o laboratório eTC
preserva uma coleta histórica de outubro de 2014, com 10,7 milhões de tweets, o
que dispensa coleta nova e torna o caso executável. Os
resultados abaixo cobrem o eixo de mídia (classificação vertical\slash
horizontal), determinístico e já rodado em escala; os eixos de stance
(preferência eleitoral) e de temas dependem de rotulagem manual ainda em
curso e não são avaliados aqui.\footnote{Números e testes deste capítulo provêm
da análise companion do projeto \texttt{pesquisa-twitter-refeita} (script
\texttt{analise/stats\_midia.py}), sobre a tabela \texttt{ituassu\_2014} do banco
do eTC. São, portanto, reprodutíveis a partir dos instantâneos versionados.}

\section{O que o artigo afirmou e o que se coletou}

O artigo analisou uma amostra de 700 tweets (100 por dia, em
horários de pico, na última semana do turno, 19--25/out), rotulados à mão,
e dela derivou, no eixo de mídia, duas afirmações testáveis: H1, o
retuíte de mídia vertical (RTMV) supera 50\% em quase toda a semana; e
H2, a mídia horizontal (MH) chega a rivalizar com a vertical (MV) em ao
menos um dia. A expansão reproduz a classificação de mídia, pelo domínio do primeiro link de
cada tweet, sobre o universo do tema da mesma janela (32.193 tweets só com \#Eleições2014;
140.909 na janela plena).

\paragraph{Ressalva metodológica.} A coleta preservada é vizinha, não idêntica,
à do artigo, por ser mais larga no escopo temático. Ela não é, portanto, um
superconjunto estrito da amostra original, e a divergência mistura, em princípio,
efeito de volume (coletar mais da mesma hashtag) com efeito de escopo (coletar
mais hashtags). Para não confundi-los, o desenho separa os dois eixos: densidade
intra-janela (E1) e largura do filtro (E2).

\section{O que se sustenta ao coletar mais (ponto positivo)}

O resultado qualitativo central da tese se sustenta. Entre o conteúdo que aponta
para mídia identificável, a vertical domina a horizontal em cerca de 5 para 1
(83,4\% de MV contra 16,6\% de MH; IC95\% de Wilson 83,1--83,6; $n=79.603$; contra
o acaso, $z=188{,}3$, $p\approx0$). Essa conclusão não depende do volume; já valia
na amostra de 700 e permanece no universo. Para ela, coletar mais não teria
alterado o resultado, e a coleta original já bastava.

\section{O que o volume maior revela (pontos negativos)}

Três afirmações quantitativas do artigo não sobrevivem à expansão. Convém adiantar,
porém, uma qualificação que a construção da curva $A(\text{volume})$ obrigou a fazer
e que a seção seguinte detalha: a expansão \emph{revelou} os erros, mas não era
necessária para \emph{evitá-los}.

\begin{description}[leftmargin=1.5em,style=nextline]
  \item[H1 (RTMV > 50\%) é artefato da amostragem por pico.] Na amostra
  reconstruída de 700, o RTMV agregado é 48,3\% (IC95\% 44,6--52,0), que não
  difere de 50\% ($p=0{,}36$). No universo da mesma janela, cai para
  36,1\% (IC95\% 35,6--36,6), decisivamente abaixo ($z=-49{,}9$,
  $p\approx0$). Os intervalos são disjuntos, sinal de viés sistemático de
  amostragem, não de flutuação. A amostra por horários de pico superestima a
  mídia vertical; no volume cheio, o conteúdo sem link (NDA, 30,8\% da janela,
  chegando a 48,8\% em 24/out) dilui a proporção. A conclusão do
  artigo reflete o recorte amostral, não o debate como um todo.
  \item[H2 (MH rivaliza com MV) não reproduz.] Em nenhum dia a MH se
  aproxima da MV; o dia que o artigo aponta como exceção (23/out) é, no universo,
  o mais vertical (MV 59\% contra MH 10,1\%). Refeita a comparação dia a dia com
  trezentas subamostras aleatórias por ponto, a mídia horizontal não alcança a
  vertical em nenhuma réplica de nenhum dia a partir de cinquenta tweets, e em
  23/out não a alcança nem com vinte e cinco. Note-se que aqui não se trata de
  volume nem de esquema: a reconstrução da própria amostra do artigo tampouco
  mostra rivalidade nesse dia (MV 62\% contra MH 13\%), e mesmo sob a leitura mais
  branda, a de que a horizontal atinge seu máximo semanal, o dia seria 21 ou
  20/out, não 23. A divergência quanto a H2 fica, portanto, sem causa atribuída.
  \item[A amostra não representa o próprio universo.] A queda de 48,3\% para
  36,1\% ($-12{,}2$ p.p., com ICs disjuntos) mostra que a subamostra do artigo
  não reconstrói nem as proporções do universo de onde saiu.
\end{description}

Num eixo adicional, o da largura do filtro (E2), abrir o escopo de \#Eleições2014
para todas as hashtags eleitorais faz a MV cair de 43,3\% para 24,9\% e o conteúdo
sem link subir para 47,1\%. A dominância medida da mídia vertical decorre, em boa
parte, de olhar apenas a hashtag-âncora. Esse deslocamento é função da hashtag,
não da composição de autores. A prova está no nível intra-autor. Os 1.752 autores
presentes nos dois recortes vão de 25,9\% para 53,3\% de conteúdo sem link (razão
de chances pareada por autor $=2{,}91$; IC95\% 2,63--3,23; Wilcoxon
$p=3{,}0\times10^{-18}$), servindo cada autor como seu próprio controle.

\paragraph{Predição refutada.} A predição registrada antes da execução supunha
que H1, por ser um efeito agregado, provavelmente não mudaria. Os dados a
refutaram. Registrar as predições antes da análise serve exatamente para expor
expectativas equivocadas como essa.

\section{A curva de volume, e por que ela corrige o diagnóstico}
\label{sec:curva-midia}

Construir a curva $A(\text{volume})$ desfaz uma atribuição de causa que os
resultados acima sugeriam. Tomando subamostras aleatórias de tamanho crescente do
universo, com trezentas réplicas por ponto, o RTMV médio é 36,1\% em \emph{todos} os
tamanhos, de cem tweets ao universo inteiro. O estimador não deriva com o volume; o
que muda é apenas a precisão, com a banda de 95\% encolhendo de dezessete pontos
percentuais, em $n=100$, para menos de um ponto em $n=22.400$.

Disso decorrem duas constatações que mudam a leitura do caso. A primeira é que os
48,3\% do artigo \emph{não cabem} na banda que uma amostra aleatória do mesmo
tamanho produz, de 32,6\% a 39,9\%: a distância não é atribuível a erro amostral, e
sim ao esquema de amostragem por horários de pico. A segunda é que bastariam
duzentos tweets sorteados ao acaso para refutar H1 com 99,7\% de poder, contra os
setecentos que o artigo de fato coletou. A conclusão correta estava, portanto, ao
alcance de uma coleta três vezes e meia \emph{menor} que a realizada.

Segue-se que dizer que ``coletar mais era necessário para chegar à conclusão
correta'' seria falso neste caso. Coletar mais revelou o erro, mas o que teria
evitado o erro era coletar de outro modo. A distinção não é escolástica: se o
diagnóstico for de volume, a recomendação prática é ampliar a coleta, que é cara; se
for de esquema, a recomendação é sortear, que é barato e frequentemente mais eficaz.

\paragraph{Uma terceira pergunta para o protocolo.} O par ``mudou?\ / convergiu?'',
como formulado no Capítulo~\ref{cap:experimento}, não distingue esses dois casos, e
aqui chegaria a atribuir ao volume uma falha de desenho amostral. Propõe-se
acrescentar, para todo artigo-alvo que amostrou em vez de analisar o universo, a
pergunta: \emph{o esquema de amostragem é não-viesado para a quantidade de
interesse?} A operacionalização é a que se usou acima, comparar a estimativa
publicada com a banda de subamostras aleatórias do mesmo tamanho. A pergunta não se
aplica quando o alvo já analisou a população, como no eixo de rede do caso do
Capítulo~\ref{cap:caso-vacinas}, mas aplica-se a todos os demais.

\section{Leitura em termos de volume}

Este caso ilustra os dois desfechos que o trabalho busca distinguir, e acrescenta um
terceiro que o desenho inicial não previa. Do lado
positivo, há uma conclusão robusta ao volume (a mídia vertical domina a
horizontal), para a qual a coleta original bastava, e que uma amostra aleatória de
cem tweets já estabeleceria. Do lado negativo, há
conclusões que a coleta pequena dava como certas e que mais dados invertem (H1,
H2 e a representatividade da própria amostra). O terceiro desfecho, que só a curva
da seção anterior tornou visível, é que essas conclusões negativas \emph{também}
haviam convergido no volume coletado: o que as inverteu não foi a falta de dados, e
sim o modo de sorteá-los. Para este caso, portanto, a pergunta ``coletar mais muda a
conclusão?'' recebe um sim que precisa ser qualificado, porque a causa não é o
tamanho da coleta.

Falta ainda o eixo de stance (a tese H3, de que os públicos diferem),
dependente de rotulagem manual. É onde o artigo faz sua afirmação mais forte e
onde o teste de volume será mais informativo. A solução para os pontos negativos,
o que fazer diante de uma conclusão dependente do volume, será proposta de forma
genérica, comum a todos os casos, e não caso a caso; o segundo caso executado
(Capítulo~\ref{cap:caso-vacinas}) fornece o contraste necessário para começar a
formulá-la.

% =====================================================================
\chapter{Segundo Caso Executado: o Debate Vacinal no Twitter}
\label{cap:caso-vacinas}
% =====================================================================

Este capítulo relata o segundo caso levado à execução: a replicação por expansão
de \citeonline{verjovsky2023quarrel}, que analisou o debate sobre vacinação no
Twitter brasileiro entre dezembro de 2021 e março de 2022. Ele entra pela mesma
exceção do caso anterior, o acervo histórico do eTC, e difere dele em dois
aspectos que o tornam complementar. Primeiro, cobre outros dois tipos de análise,
detecção de comunidades e análise de conteúdo, ausentes do eixo de mídia. Segundo,
e mais importante, aqui a sub-coleta não está na coleta, que foi larga, e sim no
\emph{estrato analisado}: a análise substantiva do artigo foi feita apenas sobre
os tweets que ultrapassaram 500 retuítes. É o exemplar mais nítido do achado do
levantamento, o de que a filtragem por critério é a norma.

Registre-se ainda que a autora desta dissertação é coautora do artigo replicado e
foi quem extraiu os dados originais. A lente ``coletar mais'' é aplicada aqui a
trabalho próprio, o que faz deste caso uma autocrítica, e não uma crítica
externa.\footnote{Números e testes deste capítulo provêm do projeto companion do
caso (scripts \texttt{analise/rede\_modularidade.py} e
\texttt{analise/curvas\_por\_limiar.py}), sobre um instantâneo congelado da coleta
preservada no eTC. São, portanto, reprodutíveis a partir do instantâneo
versionado.}

\section{O que o artigo afirmou e o que se sub-coletou}

O artigo coletou mais de 1 milhão de tweets e cerca de 4 milhões de retuítes, a
partir de seis termos-índice em português e do critério de atividade de retuíte no
período. Sobre esse corpus completo rodou a análise de rede, obtendo comunidades
de modularidade que identificou como os polos pró e anti-vacina. Já a análise de
conteúdo, conduzida pelo método de Bardin, com rotulagem manual por dois autores e
um terceiro árbitro, foi feita \emph{apenas} sobre os 1.525 tweets com mais de 500
retuítes. Dela saem o ranking de enquadramentos e a conclusão que dá título ao
trabalho, a de que os anti-vacina impuseram o enquadramento do debate.

A assimetria é o objeto do caso. A afirmação-título é, formalmente, uma afirmação
sobre a elite viral, generalizada para ``o debate no Twitter''. Como coletar mais
não é possível na plataforma, a expansão aqui é descer o limiar: reproduzem-se as
análises sobre 29.978 tweets originais em português com mais de 10 retuítes, cerca
de 19 vezes o estrato do artigo, e a rede sobre o corpus completo (879.189 autores
e 3,34 milhões de arestas de retuíte). A rotulagem em escala usa um classificador
por modelo de linguagem, validado contra os 1.525 tweets já rotulados à mão pelos
autores, que servem de gabarito.

\paragraph{Ressalva metodológica.} A coleta preservada é uma re-extração de maio de
2022, posterior ao artigo, de modo que as contagens de retuíte são maiores que as
originais e os limiares não selecionam exatamente os mesmos objetos. Além disso,
cerca de 16\% dos tweets originais não estão em português e ficam fora do eixo de
conteúdo, embora permaneçam na rede, que o artigo também rodou sem filtro de
idioma.

\section{O que se sustenta ao descer o limiar (pontos positivos)}

\begin{description}[leftmargin=1.5em,style=nextline]
  \item[A concentração é estrutural, não artefato do corte.] O artigo afirma que os
  dez maiores autores concentram 15,7\% dos retuítes. Com o denominador do artigo,
  o da amostra viral, o valor se reproduz ao decimal. No corpus completo a
  concentração cai para 9,3\%, mas permanece expressiva: dez contas respondem por
  cerca de um em cada onze retuítes de um corpus de 4,5 milhões.
  \item[A topologia confirma a assimetria de articulação entre os polos.] O
  cruzamento com os 602 autores influentes que o artigo publica é integral, com
  602 de 602 localizados na nossa partição, e revela que o lado anti-vacina forma
  \emph{uma} comunidade coesa, que reúne 244 desses autores contra um único do
  lado oposto, enquanto o lado pró se espalha por sete comunidades. A tese do
  artigo, a de um polo anti articulado contra um polo pró desarticulado, aparece
  assim na própria topologia da rede, e não apenas na leitura do conteúdo.
  \item[A composição temática convergiu.] Das 14 categorias de enquadramento, o
  maior deslocamento entre o estrato viral e o corpus 19 vezes maior é de 8,4
  pontos percentuais, e a maioria fica dentro de $\pm2$ pontos. Descer o limiar
  não reordena os temas: a tabela de prevalência temática do artigo teria sido
  essencialmente a mesma com muito mais dados.
  \item[A vantagem dos tweets pró em retuítes se mantém.] O artigo afirma que os
  tweets pró mais populares receberam o dobro de retuítes dos anti mais populares.
  Na réplica, o tweet pró mais retuitado recebe 2,33 vezes o do lado oposto, e a
  razão é a mesma em todos os limiares. Convém notar, porém, que essa estabilidade
  é estrutural, não empírica: uma afirmação sobre a cauda extrema é imune à
  sub-coleta, porque a cauda é justamente o que qualquer coleta captura primeiro.
\end{description}

\paragraph{Predição refutada.} A predição registrada antes da execução era
exatamente a inversa nesse ponto. Apostou-se que os agregados de baixa dimensão
(comunidades e concentração) se manteriam, o que ocorreu apenas em parte, e também
que o ranking de enquadramentos se recomporia abaixo do limiar, enfraquecendo a
conclusão-título por essa via. Foi a aposta interessante, e os dados a refutaram: o
que se moveu não foi o enquadramento, e sim o equilíbrio de forças, que não estava
na aposta.

\section{O que o volume maior revela (pontos negativos)}

\paragraph{O equilíbrio entre os polos não replica na rede.} A afirmação de que os
dois lados têm volumes semelhantes, cada um perto de dois quintos, é a única parte
do artigo que a expansão contradiz já no eixo de rede. Atribuído o lado a todas as
11.664 comunidades da partição pela regra do artigo, isto é, pela maioria dos autores
influentes que ele mesmo publica, a razão entre os polos fica entre 1,50 e 1,56 para
um, a favor do lado pró, contra a razão de 1,02 que a tabela do artigo reporta. O
intervalo cobre oito cenários, combinando duas bases de contagem com quatro
exigências de evidência mínima por comunidade, e a razão praticamente não se move
entre eles. O contraste fica mais nítido
em valores absolutos: o lado anti passa de 1,98 para 2,01 milhões de posts,
praticamente inalterado, enquanto o pró vai de 2,02 para 3,13 milhões.

\paragraph{A divergência, porém, não é efeito de volume.} Construída a curva
$A(\text{volume})$ da rede, sorteando frações crescentes da coleta e refazendo
partição e atribuição a cada ponto, a razão entre os polos sobe de forma monotônica
de 1,31, com 1\% dos dados, a 1,55 no corpus inteiro, e só entra na banda do valor
final a partir de três quartos da coleta. Duas leituras se seguem. A primeira é que
a conclusão \emph{não havia convergido}: com metade da coleta, o número ainda estava
em movimento, e a estrutura de comunidades converge mais devagar ainda, permanecendo
a 0,12 do teto de similaridade em três quartos do material. A segunda é uma correção
de rota: em nenhuma fração a réplica se aproxima do 1,02 relatado, de modo que a
distância para o artigo \emph{não} se explica pelo tamanho da coleta. Se explicasse,
coletar menos aproximaria os dois números, e não aproxima. O que a separa é a
convenção de atribuição de lado, e a pista está em que o artigo deixa 22,3\% dos
posts sem lado onde a réplica, sob a convenção dele, deixa 9,4\%: quanto mais
material se recusa a classificar, mais o resultado tende ao equilíbrio.

\paragraph{A mesma assimetria aparece na unidade autor, e ali o instrumento se
valida.} O artigo reporta ainda que, entre os 602 autores influentes, 58,6\% são
pró e 40,9\% anti. Reconstruindo esse conjunto pelo mesmo limiar de 500 retuítes e
atribuindo o lado \emph{pela topologia}, sem ler o texto, obtêm-se 58,9\% e 41,1\%,
uma diferença de 0,3 ponto percentual. A coincidência importa por dois motivos. O
primeiro é que valida a regra de atribuição de lado: o artigo rotulou esses autores
lendo o que escreveram, a réplica os rotulou pela comunidade de retuítes a que
pertencem, e os dois caminhos chegam ao mesmo lugar. O segundo é que, validada a
regra no ponto original, o que ela mostra fora dele passa a ter peso. E o que ela
mostra é que baixar o limiar que define ``influente'' dobra a assimetria: a razão
entre autores pró e anti vai de 1,43 acima de 500 retuítes para 2,34 acima de 25, e
não há sinal de estabilização dentro da faixa observada.

\paragraph{A cobertura agregada, em contrapartida, não serve de teste.} A outra
metade da mesma afirmação, a de que os dois polos somam 78\% dos posts, varia de
76,7\% a 90,6\% conforme a base de contagem adotada, e fica portanto sem veredito. O
contraste entre as duas metades é instrutivo: a soma dos polos depende de quantos
usuários se deixa fora, decisão de convenção que os dois trabalhos tomam de modo
diferente e nenhum dos dois documenta por inteiro; a razão entre os polos é imune a
essa escolha, porque a convenção afeta os dois lados da fração. Ao replicar por
expansão, convém identificar antes quais quantidades são invariantes à convenção,
porque só essas suportam comparação entre trabalhos.

\paragraph{O balanço entre os polos também não sobrevive no conteúdo.} Entre os tweets a
que se atribui um lado, a fatia pró-vacina cresce de forma monotônica à medida que
o limiar cai: de 48,4\% em mais de 500 retuítes para 62,6\% em mais de 10, no
esquema binário do artigo, e de 50,0\% para 61,7\% num esquema que admite classe
neutra. A direção e a magnitude do deslocamento, portanto, não dependem da escolha
de anotação. No estrato viral, o instrumento reproduz o gabarito humano publicado
com exatidão, dando os mesmos 48,4\%, o que descarta atribuir a diferença ao
classificador.

A leitura de que o polo anti-vacina tem presença equivalente ou superior vale,
assim, para o material que viralizou, não para o debate. A conclusão-título se
apoia num estrato em que ela é verdadeira e que não representa o corpus de onde
saiu.

O mecanismo aparece quando se admite uma classe neutra: a fatia de tweets
topicamente relevantes que não emitem juízo cresce de 24,6\% no estrato viral para
33,4\% no corpus amplo. A viralidade seleciona conteúdo que toma partido. É um
achado próprio do caso, invisível sob o esquema de anotação do artigo, e explica
por que um recorte por engajamento superestima sistematicamente a polarização do
debate que pretende descrever.

Convém sublinhar que os achados acima são medições independentes, em três unidades
distintas: posts, autores influentes e tweets rotulados. As duas primeiras derivam da
topologia dos retuítes, sem ler uma linha de texto; a terceira, da rotulagem do
conteúdo, sem olhar a estrutura da rede. Não compartilham instrumento, dado de
entrada nem modo de falhar, e ainda assim apontam na mesma direção. Para o protocolo,
isso sugere duas checagens que valem generalizar. Quando análises de tipos distintos
incidem sobre a mesma afirmação, a concordância entre elas é evidência mais forte do
que qualquer uma isolada. E quando uma delas reproduz um número publicado, como aqui
ocorre na unidade autor, ela serve de aferição para as demais, que em geral não têm
gabarito contra o que se medir.

\section{Volume e largura do filtro respondem ao contrário}

Além do eixo de volume, o caso permite variar a \emph{largura do filtro}, o eixo E2
do protocolo, recompondo as análises por termo-índice. O artigo coletou por seis
termos em português, e refazer as medições termo a termo produz o contraste mais
instrutivo do caso.

O balanço entre os polos \emph{não se inverte em nenhum} dos seis recortes: a razão
entre pró e anti varia de 1,41, no termo mais genérico, a 7,42 em ``anti-vax'', mas
o lado pró é majoritário em todos. É a evidência mais forte de que essa conclusão
não é artefato da lista de termos. Já a composição temática faz o oposto: o tema
dominante passa de política, com 32,8\% no termo genérico, a pessoas anti-vacina,
com 83,9\% em ``anti-vax'', e a crianças, com 88,4\% em ``vacinação infantil''.
Parte dessa variação é tautológica, pois filtrar por um termo seleciona o assunto
desse termo, e é justamente aí que está o ponto: o ranking de enquadramentos
publicado é função direta da lista de termos, embora seja apresentado como descrição
do debate. Como os termos são muito desiguais em produtividade, de 16.980 tweets a
sete, o ranking publicado é, na prática, o dos termos genéricos.

Duas observações de passagem, ambas sobre a coleta e não sobre o debate. A primeira
é que os termos destinados a capturar o polo anti-vacina capturam sobretudo o polo
oposto falando \emph{sobre} ele: em ``anti-vax'', 75,5\% dos tweets são pró-vacina e
83,9\% se enquadram na categoria ``pessoas anti-vacina''. A segunda é que um dos seis
termos, ``anti-vacinação'', devolve sete tweets em quase trinta mil; a lista efetiva
tem quatro termos, não seis.

O quadro que emerge é que volume e largura de filtro são eixos independentes e
podem inverter-se um em relação ao outro. O enquadramento é robusto ao volume e
frágil ao filtro; o balanço entre polos é o contrário. Reportar apenas um dos dois
eixos, como é praxe, dá uma segurança que os dados não sustentam.

\section{Um efeito de sub-coleta sobre o instrumento}

O caso produziu um resultado que não estava previsto e que diz respeito ao
método, não ao dado. O mesmo corpus, rotulado pelo mesmo modelo, sob dois esquemas
de anotação, admite duas descrições incompatíveis: o esquema do artigo lê, acima
de 10 retuítes, 58,8\% de pró contra 35,1\% de anti, um debate de maioria pró
clara; o esquema de três classes lê 41,1\%, 25,5\% e 33,4\% sem posição, um debate
em que um terço do material não toma partido. São cerca de 18 pontos vindos de uma
escolha de anotação, e não do dado.

A origem da divergência é rastreável e não é erro de codificação. A terceira classe
do artigo, descrita como \emph{non-relevant/ambiguous}, operou como marcador de
irrelevância tópica, e não de ausência de posição: os três casos assim
classificados no gabarito são usos da palavra ``vacina'' fora do tema. Trata-se,
portanto, de um esquema binário forçado, desenho comum na literatura de análise de
conteúdo, em que o tweet relevante mas não avaliativo é resolvido em um dos polos.

Duas consequências. A primeira é uma qualificação do artigo replicado: parte do que
suas proporções registram como adesão a um polo é conteúdo sem posição alocado por
necessidade do esquema, o que limita a generalização dessas proporções ao debate
como um todo. A segunda é metodológica e mais geral. Um classificador calibrado
contra um gabarito aprende a convenção de codificação daquele gabarito, e não a
tarefa: o classificador aqui reproduz a codificação publicada melhor
($\kappa=0{,}759$) do que uma coautora do próprio artigo recodificando o mesmo
material ($\kappa=0{,}350$), porque ela aplica um critério de neutralidade que a
convenção original não tinha. Como o gabarito só existe no estrato sub-coletado, o
instrumento importa as propriedades desse estrato ao ser aplicado ao corpus amplo.
Em rotulagem manual cega no estrato baixo, os casos neutros que o esquema binário
força para um lado se dividiram de forma aproximadamente simétrica entre os polos,
o que atenua as estimativas em direção a 50\%: os percentuais acima devem ser
lidos como limite inferior do deslocamento, e não como seu valor.

\paragraph{Confiabilidade e seu teto.} O classificador de três classes atinge
$\kappa=0{,}697$ contra codificação humana, num teto de tarefa medido em
$\kappa=0{,}746$ por reteste intracodificador. As comparações \emph{entre} estratos
são mais robustas que os valores absolutos, porque o erro do instrumento incide
igualmente nos dois lados da comparação. Registre-se que o artigo original não
reporta concordância entre codificadores, o que impede estimar quanto dessa
variação já existia na codificação publicada.

\section{Leitura em termos de volume}

Este caso é o primeiro em que os dois desfechos aparecem dentro do \emph{mesmo}
artigo-alvo e sobre o \emph{mesmo} corpus. A composição temática convergiu: o
estrato viral já representava bem a distribuição de enquadramentos, e coletar
dezenove vezes mais não teria mudado essa conclusão. O balanço entre os polos
mudou, e inverteu o sinal da leitura. Não é, portanto, que o trabalho esteja
sub-coletado ou não esteja; é que tipos de análise diferentes têm sensibilidades
diferentes à sub-coleta. Uma tabela de prevalência temática aguenta o estrato
viral; uma afirmação sobre o equilíbrio de forças do debate, não.

O caso deixa uma pendência declarada: o classificador foi validado contra gabarito
apenas no estrato viral, que é onde o gabarito existe; a rotulagem manual no estrato
baixo sustenta a direção do achado, mas com amostra pequena.

O percurso até o resultado da rede vale, ele próprio, como advertência metodológica,
em dois tempos. Primeiro, uma divergência aparente de 17 pontos na cobertura acabou
sendo comparação mal pareada: o artigo reporta como ``dois grupos'' a união dos
grupos de modularidade mapeados a cada lado, e não as duas maiores comunidades da
partição, que era o que se computava. Depois, corrigido o pareamento, a cobertura
mostrou-se instável à base de contagem, ao passo que a razão entre os polos se
manteve fixa. Uma divergência grande demais para ser real costuma ser sintoma de
unidades trocadas; e uma concordância boa demais, obtida numa única convenção, pede
ser testada em outra antes de virar resultado.

Registre-se, por fim, o que não se sabe. A regra de atribuição, fiel à do artigo, faz
uma comunidade difusa de 349 mil autores herdar um lado por inteiro; é uma
generosidade que os dois trabalhos compartilham. A coleta preservada tem 12\% mais
retuítes que a do artigo dentro da mesma janela, resíduo da re-extração que segue sem
explicação. E o artigo não publica a distribuição por grupo, apenas o agregado, o que
impede conferir a decomposição dele contra a nossa. O que se afirma aqui, portanto, é
que o equilíbrio relatado não se reproduz, não que se conheça a razão completa disso.

% =====================================================================
\chapter{Terceiro Caso Executado: Idosos e Pandemia no YouTube}
\label{cap:caso-youtube}
% =====================================================================

Este capítulo relata o terceiro caso levado à execução, e o primeiro fora do
Twitter/X: a replicação por expansão de~\citeonline{ghenai2025agecovp}, que analisou
o discurso sobre pessoas idosas durante a pandemia de COVID-19 no YouTube. Ele
acrescenta ao conjunto uma rede nova e um tipo de análise que ainda não havia sido
coberto por caso executado, a modelagem de tópicos.

A escolha do alvo foi determinada pela \emph{rota de coleta}, não pela citação, e o
critério merece registro porque é, ele mesmo, um achado. Das três redes ainda
pendentes no conjunto de casos, duas se mostraram inviáveis por motivos externos ao
mérito científico: o acesso à Ad Library da Meta exige conta com documento
verificado, e a credencial de pesquisa do TikTok disponível ao grupo havia sido
revogada silenciosamente entre junho de 2025 e agosto de 2026. A API de dados do
YouTube, ao contrário, permanece gratuita, é obtida em minutos e não exige
verificação. Além disso, o objeto persiste: vídeos e comentários publicados em 2020
continuam recuperáveis em 2026, o que distingue este caso de estudos sobre
\emph{recomendação}, cujo instrumento é o próprio algoritmo e não sobrevive à
passagem do tempo.

\paragraph{Ressalva de tração.} O artigo replicado foi publicado em agosto de 2025 e
ainda não acumulou citações. Não se trata de escolha descuidada: entre os dezoito
trabalhos de YouTube do levantamento que declaram amostragem, todos são de 2025 ou
2026, e o mais citado reúne três citações. Análise computacional de YouTube com
sub-coleta declarada e tração bibliométrica não coexistem hoje. O critério que
restou foi o veículo, e a \emph{EPJ Data Science} é periódico revisado por
pares.\footnote{Números deste capítulo provêm dos scripts
\texttt{analise/fase2\_ponto\_original.py}, \texttt{analise/fase3\_efeito\_filtro.py}
e \texttt{analise/fase3\_no\_tema.py} do projeto companion do caso, sobre o conjunto
de dados publicado pelos autores e sobre um instantâneo próprio de coleta.}

\section{O que o artigo fez e o que se sub-coletou}

Os autores cruzaram dezessete termos relativos a pessoas idosas com cinco relativos à
pandemia, totalizando oitenta e cinco buscas, e recolheram até seiscentos resultados
de cada uma. Desse ponto em diante, a coleta é sucessivamente filtrada:

\begin{itemize}
  \item das 6.997 respostas iniciais, mantêm-se as 3.353 em que a palavra buscada
  \emph{reaparece no título ou na descrição} do vídeo, descartando 52\%;
  \item para cada vídeo mantido, coletam-se os vídeos sugeridos pela plataforma,
  104.172 ao todo, dos quais o mesmo critério de palavra-chave preserva 1.025,
  descartando 99\%;
  \item a união passa por um filtro de idioma e resulta nos 3.782 vídeos finais, com
  2.243 canais e 69.425 comentários.
\end{itemize}

O critério que governa as duas etapas é de \emph{forma}, não de assunto: exige que o
texto do vídeo repita a palavra da busca. É filtragem por critério no sentido exato
do levantamento, e a hipótese que organiza este caso é que ela não seleciona ao
acaso. Quem redige título repetindo a palavra-chave é, tipicamente, veículo de
imprensa; quem publica um relato pessoal sobre a avó doente não escreve assim.

A expansão, aqui, é soltar esse filtro. Recolheram-se as mesmas buscas guardando
\emph{todos} os resultados e registrando a aprovação no filtro como atributo, e não
como critério de exclusão, de modo que o mesmo instantâneo sirva aos dois lados da
comparação.

\section{O ponto de partida: o que se reproduz}

Antes de expandir, é preciso saber reproduzir. Os autores publicaram o conjunto
completo, o que permitiu verificar o ponto original sem consumir coleta nova. Das
afirmações verificáveis, duas se reproduzem e uma se reproduz pela metade.

A composição por categoria de canal reproduz a menos de dois pontos percentuais,
mas somente sob uma convenção que o artigo não declara: contar apenas a primeira
categoria de cada canal. As demais convenções plausíveis, contar cada canal em todas
as suas categorias ou contar atribuições, erram entre vinte e sete e setenta e cinco
pontos percentuais. A afirmação comparativa central, a de que os comentários são mais
negativos que os vídeos, reproduz com folga larga: a toxicidade média dos comentários
é setenta e quatro vezes a dos vídeos. Já a distribuição de sentimento reproduz em
uma das duas bibliotecas empregadas e não na outra, e a distância na segunda depende
de um limiar de neutralidade que o artigo também não publica.

O padrão que emerge das três verificações é uniforme e vale registrar: onde houve
divergência, a causa foi sempre a mesma, um parâmetro de convenção deixado implícito.

\section{Uma afirmação sem valor de verdade determinado}

Uma quarta afirmação do artigo não pôde ser testada, e o motivo é de outra natureza.
O texto sustenta que o corpus é dominado por veículos de imprensa e não por conteúdo
de usuário comum, e usa essa leitura para explicar dois resultados distintos, o
sentimento positivo dos vídeos e sua baixa toxicidade. Ao procurar a regra que
separa as duas categorias, encontrou-se não a regra, mas uma contradição interna.

Ao relatar a rotulagem manual, o artigo afirma que \emph{menos de 7\% dos vídeos
eram conteúdo copiado, e não conteúdo de usuário}, o que implica que mais de 93\%
seriam de usuário. A frase imediatamente seguinte afirma o contrário, que o conjunto
é dominado por vídeos originalmente veiculados na televisão e que pouquíssimos foram
criados por usuários. Adiante, na seção de sentimento, o mesmo número reaparece
invertido, agora atribuindo aos usuários a fatia inferior a 7\%. As estatísticas de
canal, que registram predomínio de canais pequenos, apoiam a primeira leitura; a
seção de toxicidade, que atribui a moderação da linguagem aos veículos, apoia a
segunda.

Os rótulos que decidiriam a questão foram produzidos manualmente, vídeo a vídeo, e
não integram o conjunto publicado. A afirmação, portanto, não é adjudicável de fora.
Uma medição independente, conduzida aqui com regra declarada, situa o conteúdo de
usuário em 32\% do corpus, valor incompatível com qualquer das duas versões
publicadas.

O caso não é isolado. No debate vacinal, o corpo do artigo descreve onze categorias
de enquadramento enquanto o material suplementar numera catorze. Nos anúncios de
imigração, o total de impressões relatado não se recompõe a partir da planilha que o
próprio artigo publica. Em três redes e três tipos de análise, o mesmo fenômeno: o
número publicado depende de uma decisão que não foi publicada. É um achado sobre
reprodutibilidade que não depende de coletar nada, e que a lente ``coletar mais''
tornou visível por acidente, ao exigir que cada afirmação fosse operacionalizada
antes de ser testada.

\section{O efeito do filtro}

A comparação ingênua entre o que o filtro aprova e o que ele descarta é enganosa, e
a primeira rodada deste caso caiu nessa armadilha. Entre os vídeos descartados
aparecem categorias como esporte e beisebol, ausentes do lado aprovado: são
respostas de busca que simplesmente não tratam do assunto. Comparar os dois conjuntos
sem separá-las mede ruído, não viés, e a predição registrada antes da medição, a de
que a fatia de canais institucionais cairia ao soltar o filtro, foi refutada por esse
motivo espúrio.

A comparação foi então refeita apenas entre vídeos que passam por um teste de tema
independente, exigindo menção a pessoas idosas \emph{e} à pandemia. O teste foi
validado sobre o próprio corpus do artigo, que ele aprova em 99,9\%. Restrita ao
tema, a medição separa dois efeitos que estavam sobrepostos:

\begin{description}[leftmargin=1.5em,style=nextline]
  \item[O filtro é, em boa medida, limpeza legítima.] Ele descarta 91,4\% dos vídeos
  fora do tema e apenas 25,0\% dos que estão no tema. A objeção genérica de que
  descarta demais não se sustenta.
  \item[Mas o resíduo é substantivo.] Dos vídeos descartados, 19,5\% estão no tema.
  Um em cada cinco descartes é perda de conteúdo pertinente.
  \item[E o descarte é enviesado na direção prevista.] Entre os vídeos no tema, e sob
  uma mesma regra aplicada aos dois lados, o conteúdo de usuário responde por 30,9\%
  do que o filtro aprova e por 38,8\% do que ele descarta, diferença de 7,9 pontos
  percentuais. O filtro remove sistematicamente mais conteúdo de usuário comum do que
  preserva.
\end{description}

A consequência substantiva desse último ponto merece desenvolvimento, por ser a
afirmação mais forte que este capítulo faz sobre o trabalho replicado.

O artigo encontra dois resultados que ele próprio classifica como inesperados para o
tema, e explica ambos pela mesma causa. O primeiro é o sentimento dos vídeos, que
apura predominantemente positivo, 59\% pela biblioteca VADER e 72\% pela TextBlob, e
que o texto atribui, entre outros fatores, à linguagem neutra de vídeos que ``em sua
maioria vêm de veículos de imprensa, e não de conteúdo gerado por usuários''. O
segundo é a toxicidade, que apura mínima, com mais de 90\% dos vídeos pontuando 0,3 ou
menos, e que o texto atribui ao fato de que ``a maior parte dos vídeos vem de mídia
noticiosa e de veículos informativos, que tipicamente empregam linguagem moderada e
neutra''. Em ambos os casos, portanto, a composição do corpus opera como causa
explicativa, e é tratada como propriedade do que se encontra na plataforma quando se
observa esse assunto.

Ocorre que essa composição é, em parte, produto da regra de seleção do próprio estudo.
Exigir que a palavra buscada reapareça no título ou na descrição privilegia uma
prática de redação editorial: um veículo intitula ``Coronavirus: elderly care homes
face crisis'', ao passo que quem publica o relato de uma avó adoecida intitula ``my
grandma's story''. O critério, sem que essa fosse a intenção, tem mais chance de
aprovar o primeiro. A medição da seção anterior quantifica esse efeito: restrita aos
vídeos que tratam do assunto, e portanto livre do ruído que contaminaria a comparação,
a fatia de conteúdo de usuário é de 30,9\% entre os vídeos que o filtro aprova e de
38,8\% entre os que ele descarta. O material descartado é sistematicamente mais
autoral.

Segue-se que atribuir a moderação da linguagem à natureza do corpus é atribuí-la, em
alguma medida, à própria decisão de coleta. Não se trata de erro de execução, e sim de
um efeito que os autores não tinham como observar, uma vez que o material descartado
não foi caracterizado. Mantém-se, porém, a qualificação: a circularidade é
\emph{parcial}, e por três razões. O filtro não é a única causa da predominância
institucional, pois o próprio ordenamento de busca do YouTube favorece canais
estabelecidos e veículos de imprensa efetivamente produziram grande volume de conteúdo
sobre pandemia e envelhecimento; a diferença medida, de 7,9 pontos percentuais, indica
direção com segurança mas não sustenta a afirmação de que o corpus se inverteria sem o
filtro; e a caracterização de canais alcançou 165 canais do lado descartado contra 531
do lado aprovado, o que torna a magnitude imprecisa. O que a evidência sustenta é que
parte da propriedade usada como explicação foi produzida pelo procedimento que a
mediu.

Registre-se que a segunda predição, sobre a composição por categoria de canal,
falhou também no recorte temático, e permanece sem causa atribuída.

\section{Limites deste caso}

Três limites devem acompanhar qualquer leitura dos números acima. O primeiro é de
completude: a coleta própria cobre cinquenta das oitenta e cinco buscas e duas das
doze páginas por busca, o que torna as proporções mais confiáveis que as contagens
absolutas. O segundo é de precisão: do lado descartado, apenas 165 canais puderam ser
caracterizados, contra 531 do lado aprovado, de modo que a diferença de 7,9 pontos é
direcionalmente confiável e numericamente imprecisa.

O terceiro é de natureza diferente e merece ênfase. O ramo mais espetacular do funil,
o dos vídeos sugeridos, com seus 99\% de descarte, \emph{não é reproduzível}: o
parâmetro da API que devolvia vídeos relacionados foi removido em agosto de 2023.
Um trabalho publicado em 2025 possui, portanto, uma etapa de coleta que já não pode
ser refeita. A janela de replicabilidade de um estudo de rede social não é apenas
estreita; ela se fecha, e às vezes antes mesmo da publicação.

% =====================================================================
\chapter{Um Quarto Caso, Parcial: Deleção de Conteúdo Político no TikTok}
\label{cap:caso-tiktok}
% =====================================================================

Este capítulo relata um caso de natureza diferente dos três anteriores. Ele não
executa o eixo ``coletar mais'', porque a coleta em TikTok não estava disponível, e
ainda assim produz dois resultados: a reprodução do ponto original de um trabalho
publicado e a identificação de um eixo de variação que o protocolo desta dissertação
não previa. Entra no conjunto por isso, e com o rótulo de parcial.

\section{Como a célula foi preenchida sem coleta}

A rota de coleta do TikTok estava fechada. A API de pesquisa da plataforma exige
credencial concedida por projeto e prazo determinado; a credencial disponível ao
laboratório havia sido revogada em algum ponto entre junho de 2025 e agosto de 2026,
sem aviso, e o coletor institucional seguia em execução sem coletar nada. Pedir
credencial nova leva semanas.

O impasse foi contornado invertendo a pergunta. Em vez de perguntar como coletar do
TikTok, perguntou-se que trabalho sobre TikTok publica dados suficientes para ser
replicado \emph{sem} coleta. A resposta foi
\citeonline{ruiz2025politok}, que reúne 938.961 publicações de duas eleições alemãs, a
estadual da Saxônia em 2024 e a federal em 2025, e disponibiliza publicamente, sob
licença aberta, dois artefatos: o identificador de cada publicação acompanhado de seu
\emph{estado de disponibilidade em datas sucessivas de verificação}, e as anotações
humanas de um subconjunto. Os autores não liberam mídia nem texto, mas o que liberam
basta para verificar as afirmações centrais, todas relativas a quanto do corpus
desapareceu da plataforma.

\section{O que se reproduz}

Das sete afirmações quantitativas do artigo, seis se reproduzem, quatro delas de
forma exata. Reproduzem-se o volume total, as taxas de deleção das duas coletas
(18,7\% na estadual e 39,7\% na federal), a proporção de deleções atribuíveis à
retirada pelo próprio autor, aproximadamente dois terços, e as duas afirmações
derivadas da anotação humana, a de que cerca de um em cada cinco posts veicula
intolerância e a de que a maioria veicula humor.

A sétima não fecha, e falha pelo motivo que já se tornou recorrente neste conjunto de
casos. O artigo afirma que a plataforma, e não o autor, removeu 13,0\% de todas as
publicações, mas não informa quais códigos de estado atribui à plataforma. Testadas
cinco convenções plausíveis sobre os códigos disponíveis, a mais próxima produz
11,6\%, a 1,4 ponto percentual do publicado. A afirmação replica aproximadamente, sob
convenção que precisou ser reconstruída.

Também as duas afirmações de anotação replicam apenas sob uma convenção não declarada:
os valores publicados correspondem à contagem \emph{por anotação}, e não por post.
Como 300 dos 360 posts anotados receberam mais de uma anotação, contar por post com
voto majoritário desloca a intolerância de 20,5\% para 16,1\%. A conclusão qualitativa
sobrevive nas duas convenções; o número, não.

\section{O eixo que o protocolo não previa}

O achado próprio deste caso vem de um artefato que o conjunto de dados oferece e que
nenhum dos casos anteriores possuía: a coleta da Saxônia foi verificada em três datas
distintas, sobre exatamente as mesmas publicações.

\begin{table}[htbp]
\centering
\caption{Taxa de deleção do mesmo conjunto de publicações, por data de verificação}
\label{tab:tiktok-temporal}
\begin{tabular}{lrr}
\toprule
Data da verificação & Dias após a eleição & Deletados \\
\midrule
02/out/2024 & 31  & 6,3\% \\
10/dez/2024 & 100 & 17,3\% \\
13/jan/2025 & 134 & 18,7\% \\
\bottomrule
\end{tabular}
\end{table}

Os valores da Tabela~\ref{tab:tiktok-temporal} são calculados sobre as publicações
efetivamente reverificadas em cada data. A ressalva importa: nas duas primeiras
verificações, cerca de 46\% do conjunto não foi reverificado, e tomar o total como
denominador subestimaria a taxa, produzindo a série 3,3\%, 9,5\% e 18,7\%.

A leitura é direta e não depende da escolha de denominador. O denominador é o mesmo
nas três datas, e nenhuma publicação foi acrescentada entre elas; ainda assim a
quantidade medida triplica em pouco mais de três meses. Não se trata de coletar mais,
mas de \emph{coletar depois}. Dois trabalhos que aplicassem o mesmo procedimento ao
mesmo corpus, um em outubro e outro em janeiro, publicariam 6,3\% e 18,7\%, e ambos
estariam corretos.

Mais interessante que a magnitude é a forma. O salto ocorre entre a primeira e a
segunda verificação; entre a segunda e a terceira o movimento é de 1,4 ponto. A
medida sobe rápido e assenta, que é exatamente o comportamento das curvas de volume
construídas nos casos anteriores. O eixo temporal tem, portanto, sua própria curva de
convergência, e admite a mesma pergunta que o protocolo já faz ao volume: a partir de
quando a medida estabiliza.

Isso amplia o quadro que os casos anteriores vinham desenhando. O debate vacinal
mostrou que volume e largura de filtro são eixos independentes, que podem responder em
direções opostas. O caso do YouTube acrescentou a forma do critério de seleção. Este
acrescenta o momento. São quatro decisões de coleta que governam o resultado
publicado, e apenas a primeira é rotineiramente reportada. Para conclusões sobre
conteúdo que a plataforma modera, a data da coleta não é detalhe de procedimento: é
parte da definição da quantidade medida.

\section{O que este caso não entrega}

Não entrega o eixo central do protocolo. Sem acesso de coleta, não há expansão a
executar, e as colunas ``mudou?'' e ``convergiu?'' da tabela consolidada permanecem
vazias no que diz respeito ao volume. O que o caso oferece no lugar é a reprodução do
ponto original e o eixo temporal, ambos obtidos exclusivamente a partir de material
publicado.

Aplica-se aqui a mesma ressalva de tração já declarada para o caso do YouTube: o
trabalho replicado é de 2025 e ainda não acumulou citações. Vale, ademais, o registro
de que a viabilidade deste caso decorreu inteiramente de uma decisão dos autores
originais, a de publicar identificadores e estados de disponibilidade sob licença
aberta. Sem essa decisão, a célula permaneceria vazia, e não por falta de mérito do
alvo.

% =====================================================================
\chapter{Conclusão Parcial e Próximos Passos}
\label{cap:conclusao}
% =====================================================================

Esta dissertação articula um levantamento sistemático companion de~\citeonline{barreto2026survey} e um protocolo que dele decorre,
tomando as RSD como caso de estudo de uma pergunta metodológica mais
ampla, a de saber se coletar pouco ou muito muda a conclusão a que uma análise chega.
O levantamento mostra que a literatura de RSD dimensiona sua coleta por
conveniência e quase nunca verifica que o volume coletado bastava; o protocolo
propõe medir, diretamente, se coletar mais muda a conclusão e se coletar menos
teria bastado, por tipo de análise e por rede social, com um conjunto fixado de
seis casos replicáveis. Três deles já foram executados. A replicação de
\#Eleições2014 no Twitter (Capítulo~\ref{cap:caso-twitter}) mostra, no eixo de
mídia, os dois lados esperados. De um lado, uma conclusão robusta ao volume, a de
que a mídia vertical domina a horizontal, que a coleta original já bastava para
estabelecer. De outro, conclusões que o volume maior revela frágeis, como o
limiar RTMV > 50\%, que era artefato da amostragem por pico. Construir a curva de
volume desse caso produziu, ainda, um resultado sobre o próprio protocolo: a
conclusão frágil \emph{já havia convergido} no volume que o artigo coletou, de modo
que o par ``mudou?\ / convergiu?'' atribuiria ao volume uma falha que é de esquema
de amostragem. Daí a terceira pergunta proposta na
Seção~\ref{sec:curva-midia}, a de se o esquema é não-viesado para a quantidade de
interesse, que passa a integrar o protocolo para todo alvo que amostra. A replicação do
debate vacinal de 2021--22 (Capítulo~\ref{cap:caso-vacinas}) reencontra os dois
lados dentro de um único artigo-alvo e sobre o mesmo corpus: a composição temática
já havia convergido no estrato analisado, ao passo que o balanço entre os polos
pró e anti-vacina, do qual depende a conclusão-título do artigo, se inverte assim
que se desce o limiar de engajamento que definiu o estrato. Os dois casos apontam,
por caminhos distintos, para a mesma leitura: a sensibilidade ao volume é
propriedade do tipo de análise, não do trabalho. A resposta obtida nas
RSD, por serem um domínio em que o efeito é mensurável, informa a questão para
além delas.

O terceiro caso executado, sobre o discurso a respeito de pessoas idosas no
YouTube (Capítulo~\ref{cap:caso-youtube}), acrescenta a essa leitura um eixo que os
dois primeiros não isolavam. Ali a sub-coleta não é de volume nem de estrato de
engajamento, mas de \emph{forma}: um filtro que exige a repetição da palavra buscada
no título ou na descrição do vídeo. Medido apenas entre os vídeos que tratam do
assunto, o filtro se revela em parte limpeza legítima, pois descarta 91,4\% do que
está fora do tema contra 25,0\% do que está dentro, e em parte viés sistemático, pois
o material que descarta é 7,9 pontos percentuais mais rico em conteúdo de usuário
comum do que o material que preserva. Como o artigo replicado atribui à predominância
de veículos de imprensa a explicação de dois de seus resultados, e como parte dessa
predominância é produzida pelo próprio filtro, a explicação é parcialmente circular.

Esse caso trouxe ainda um achado de natureza distinta, que atravessa os três e não
depende de coletar coisa alguma. Em cada um dos alvos replicados, ao menos um número
publicado revelou-se dependente de uma decisão que não foi publicada: a convenção de
contagem de categorias e o limiar de neutralidade no caso do YouTube, a numeração das
categorias de enquadramento no do debate vacinal, a regra de agregação de impressões
no dos anúncios de imigração. No limite, uma afirmação chegou a receber, no mesmo
artigo, dois valores mutuamente exclusivos. A lente ``coletar mais'' tornou isso
visível por um efeito colateral do método: para testar se uma conclusão muda, é
preciso primeiro operacionalizá-la, e é nesse momento que a decisão ausente aparece.
Trata-se de um resultado sobre reprodutibilidade que o protocolo produz de graça, ao
lado do resultado sobre volume que ele busca deliberadamente.

Um quarto caso, deliberadamente parcial, fecha o quadro por outro lado
(Capítulo~\ref{cap:caso-tiktok}). Sem acesso de coleta ao TikTok, ele não executa a
expansão; reproduz o ponto original de um trabalho sobre deleção de conteúdo político
e, ao fazê-lo, expõe um eixo que este protocolo não previa. Verificado três vezes ao
longo de três meses, o mesmo conjunto de publicações apresenta taxas de deleção de
6,3\%, 17,3\% e 18,7\%: a quantidade medida triplica sem que uma única publicação
tenha sido acrescentada ao corpus. O movimento tem, ademais, a forma já familiar das
curvas de volume, subindo rápido e assentando, o que sugere que também o momento da
coleta admite a pergunta sobre convergência.

Somados, os casos executados identificam quatro decisões de coleta capazes de governar
a conclusão publicada: o \emph{volume} recolhido, a \emph{largura} do filtro temático,
a \emph{forma} do critério de seleção e o \emph{momento} da verificação. Apenas a
primeira é reportada com regularidade na literatura levantada, e é também a única que o
protocolo, tal como formulado no Capítulo~\ref{cap:experimento}, mede diretamente.
Estendê-lo às outras três é o desdobramento natural deste trabalho, e a evidência
reunida aqui sugere que a extensão não é ornamental: nos casos em que dois eixos
puderam ser medidos sobre o mesmo corpus, eles responderam de modo independente, e por
vezes em direções opostas.

\section{Limitações conhecidas}
\label{sec:limitacoes}

\begin{enumerate}
  \item Acessibilidade das plataformas: o protocolo depende de coleta
  ainda viável hoje. O Twitter/X, central na literatura, só entra no
  conjunto de casos porque o eTC preserva um acervo histórico de 2014
  (Capítulo~\ref{cap:caso-twitter}); é uma exceção que não se reproduz por coleta
  nova, e o CrowdTangle, sem acervo equivalente, fica de fora. A dependência
  de acervos pré-existentes é, ela mesma, uma limitação de generalização.
  \item Superconjunto aproximado: no caso do Twitter, o acervo disponível
  é uma coleta vizinha, não um superconjunto estrito da amostra do artigo;
  o efeito de volume é, por isso, separado do efeito de escopo de coleta em vez de
  confundido com ele (Capítulo~\ref{cap:caso-twitter}).
  \item Custo da expansão: coletar superconjuntos (especialmente sob a
  cota semanal da Meta Content Library) impõe um teto prático ao $N_{\text{máx}}$
  por caso.
  \item Confusão coleta versus método: mitigada fixando o modelo/
  algoritmo de cada análise ao do artigo original, isolando o efeito da coleta.
  O segundo caso executado mostra que essa fixação não basta na análise de
  conteúdo: o esquema de anotação é ele próprio uma escolha de método, e dois
  esquemas aplicados ao mesmo corpus pelo mesmo modelo produzem descrições que
  diferem em cerca de 18 pontos percentuais, ainda que a comparação entre estratos
  permaneça robusta a essa escolha (Capítulo~\ref{cap:caso-vacinas}). Some-se que
  um classificador calibrado contra gabarito aprende a convenção de codificação
  nele registrada, e o gabarito, por construção, só existe no estrato sub-coletado.
  \item Explicação (o ``porquê'') fora do escopo: reporta-se se e
  para quais tipos de análise o volume de coleta muda a conclusão, não
  por que uma análise é sensível ou robusta ao volume. O mecanismo (por exemplo, por que a análise de sentimento reagiria a um conjunto
  maior enquanto a estatística descritiva não) fica como agenda, não como
  resultado.
  \item Teoria de amostragem fora do escopo: a questão estatística de
  sampling (garantias de representatividade, inferência sobre um todo) não
  é tratada; o efeito do volume é medido empiricamente, comparando coletar
  mais e menos sobre o mesmo tema, sem modelar a coleta como esquema de
  amostragem inferencial.
\end{enumerate}

\section{Próximos passos}

\begin{itemize}
  \item fechar o caso do Twitter: rotular stance (EA/ED) e temas,
  validar os classificadores contra rótulo humano e testar a tese central H3 (os
  públicos diferem?), o eixo em que o artigo faz sua afirmação mais forte;
  \item fechar o caso do debate vacinal: ampliar a
  validação do classificador fora do estrato viral, onde o gabarito publicado não
  alcança, e investigar se o crescimento do corpus concentrado no polo pró decorre
  da janela de re-coleta ou do próprio fenômeno;
  \item executar os demais casos do conjunto de casos (Reddit, YouTube, TikTok e
  Instagram/Facebook), começando pelos de centralidade e detecção de comunidades;
  \item consolidar a tabela mestre e as curvas $A(\text{volume})$ por caso;
  \item com dois casos em mãos, e mais tipos de análise à medida que os demais
  forem executados, propor uma solução genérica para os pontos negativos (o que
  fazer diante de uma conclusão dependente do volume), comum a todos os casos, em
  vez de uma correção caso a caso.
\end{itemize}

% =====================================================================

  \arial
  \bibliographystyle{abnt-alf}
  \bibliography{referencias}
\end{document}
